\chapter{Considerações Finais} \label{cap:consideracoes}

%\section{Conclusões do Projeto de Formatura}
%Apresentar o balanço do trabalho: resultados atingidos e não atingidos, com justificativas.

\section{Contribuições}
Além dos resultados práticos do protótipo, este trabalho apresenta três contribuições
técnicas de caráter mais amplo à comunidade de Engenharia:
 
\subsection{Reprodutibilidade por meio do ESP32}
 
O \textit{dongle} é construído sobre o microcontrolador \textbf{ESP32}, um
componente amplamente disponível, de baixo custo e bem documentado. Essa escolha
garante que qualquer pesquisador, estudante ou profissional possa reproduzir,
adaptar ou estender o hardware sem dependência de componentes proprietários ou de
difícil aquisição. Todo o esquemático, \textit{layout} de PCB e código-fonte do
firmware serão disponibilizados publicamente, tornando o trabalho integralmente
reprodutível.
 
\subsection{Entregável \textit{Open Source} como Base para Pesquisas Futuras}
 
O \textit{dongle} desenvolvido será publicado como projeto \textit{open source}
(hardware e software), servindo como plataforma de referência para pesquisas que
demandem coleta de dados veiculares confiáveis. Trabalhos nas áreas de
\textit{driver behavior}, manutenção preditiva, controle de frotas e detecção de
anomalias em barramentos CAN poderão utilizar o entregável como ponto de partida,
sem precisar lidar com as vulnerabilidades presentes nos \textit{dongles} comerciais
atualmente disponíveis.
 
\subsection{Protocolo Padronizado e Seguro de Comunicação OBD-II}
 
A principal lacuna identificada na literatura é a ausência de um protocolo
padronizado de comunicação entre \textit{dongles} OBD-II e \textit{companion apps}
que incorpore requisitos de segurança desde a sua concepção. Este trabalho irá implementar um protocolo baseado no existente ELM-327 \cite{elm327datasheet}, porém com considerações de segurança que atendam os serviços de integridade, confidencialidade e autenticidade.

 
%\section{Perspectivas de Continuidade}
%Descrever os trabalhos que podem ser realizados como continuação do projeto de formatura.